% Informe de cierre Sprint 3 - CuscoNodes
\documentclass[12pt,a4paper]{article}
\usepackage[utf8]{inputenc}
\usepackage[T1]{fontenc}
\usepackage[spanish]{babel}
\usepackage{graphicx}
\usepackage{hyperref}
\usepackage{array}
\usepackage{longtable}
\usepackage{caption}
\usepackage{geometry}
\usepackage{setspace}
\usepackage{booktabs}
\geometry{left=2.5cm,right=2.5cm,top=2.5cm,bottom=2.5cm}
\setlength{\parskip}{0.5em}
\setlength{\parindent}{0pt}
\renewcommand{\familydefault}{\rmdefault}
\usepackage{xcolor}
\color{black}
\newcommand{\sectiontitle}[1]{\vspace{0.6cm}\noindent{\fontsize{13}{15}\selectfont\textbf{#1}}\par\vspace{0.15cm}\noindent{\fontsize{12}{14}\selectfont}}
\newcommand{\bodytext}[1]{{\fontsize{12}{14}\selectfont #1}\par}
\begin{document}

% Portada
\begin{titlepage}
  \centering
  {\bfseries\Large UNIVERSIDAD ANDINA DEL CUSCO}\\[0.5cm]
  {\large FACULTAD DE INGENIERÍA Y ARQUITECTURA}\\[0.2cm]
  {\large ESCUELA PROFESIONAL DE INGENIERÍA DE SISTEMAS}\\[0.8cm]

  % Logo (opcional)
  \includegraphics[width=0.35\textwidth]{logo_uac.png}\\[0.6cm]

  {\bfseries\Large INFORME DE CIERRE DE SPRINT}\\[0.4cm]
  {\bfseries\Large CuscoNodes - Sprint 3: Motor de Razonamiento}\\[1.0cm]

  \begin{flushleft}
    \textbf{ASIGNATURA:}\hspace{1.5cm} Desarrollo de Proyecto / Ingeniería de Sistemas\\[0.3cm]
    \textbf{DOCENTE:}\hspace{2.3cm} Mg. Ing. Hugo Espetia Huamanga\\[0.3cm]
    \textbf{GRUPO:}\hspace{2.8cm} HABA\\[0.8cm]
    \textbf{INTEGRANTES:}\\[0.2cm]
    \hspace{3cm} Morales Ttito Honorio \hfill 100\\%
    \hspace{3cm} Torres Cáceres Sayo Michael \hfill 100\\%
    \hspace{3cm} Zavaleta Luna Sebastian (SM) \hfill 100\\%\\
  \end{flushleft}
  \vspace{1.5cm}

  \href{https://github.com/Honorio-Morales/cusconodes}{https://github.com/Honorio-Morales/cusconodes}\\[0.8cm]

  \vfill
  {\large CUSCO -- PERÚ}\\[0.3cm]
  {\large 2026}
\end{titlepage}

% Índice
{\fontsize{12}{14}\selectfont\tableofcontents}\newpage

% Introducción
\sectiontitle{Introducción}
\bodytext{Este documento presenta el cierre formal del Sprint 3 del proyecto \textit{CuscoNodes}. Describe objetivos, trabajo realizado, resultados técnicos, evidencias y recomendaciones para la transición al Sprint 4 (Acción: traducción y notificaciones).}

% Objetivos
\sectiontitle{Objetivos}
\bodytext{El objetivo principal del Sprint 3 fue diseñar, implementar e integrar un Motor de Razonamiento local que clasifique alertas por urgencia y genere recomendaciones, evitando dependencias externas pesadas.}
\begin{itemize}
  \item Implementar un clasificador local (sklearn) para urgencia (CRÍTICA/INFORMATIVA/IRRELEVANTE).
  \item Detectar ubicaciones turísticas y tipos de eventos relevantes para Cusco.
  \item Integrar el Motor de Razonamiento en el pipeline existente (Percepción → Razonamiento → Almacenamiento).
  \item Proveer una salida JSON consumible por la interfaz web y por futuros módulos de traducción/notificación.
\end{itemize}

% Desarrollo del Sprint
\sectiontitle{Desarrollo del Sprint}
\bodytext{Se optó por una implementación "lite" con scikit-learn (TfidfVectorizer + MultinomialNB) para asegurar ejecución local, rapidez y reproducibilidad. El agente fue entrenado inicialmente con ejemplos sintéticos contextualizados al dominio turístico de Cusco y evaluado con 1193 alertas reales de Sprint 2.}

\subsection*{Arquitectura implementada}
\bodytext{El Motor de Razonamiento se ubica como segunda etapa del pipeline:}
\begin{itemize}
  \item Etapa 1 - Percepción: Scrapers (Sprint 2)
  \item Etapa 2 - Razonamiento: `ReasoningAgentLite` (Sprint 3)
  \item Etapa 3 - Acción: (Sprint 4 planificado)
\end{itemize}

% Resultados
\sectiontitle{Resultados}
\bodytext{Resumen de resultados obtenidos durante validaciones y pruebas locales.}
\begin{longtable}{p{6cm} p{8cm}}
\toprule
\textbf{Entrega} & \textbf{Evidencia / Notas} \\
\midrule
Agente de Razonamiento Ligero & `src/reasoning_agent_lite.py` (Tfidf + MultinomialNB) \\
Integración en Scheduler & `src/scheduler.py` actualizado; salida en `/data/processed/alertas_clasificadas_*.json` \\
Servidor demo web local & `web_server.py` exponiendo `/api/alerts/latest` \\
Demostración con datos reales & `sprint3_reasoning_demo.py` procesó 1193 alertas en <1s \\
Documentación de cierre & `docs/SPRINT3_FINAL.md`, `docs/INFORME_SPRINT3.tex` \\
\bottomrule
\end{longtable}

\subsection*{Estadísticas de clasificación}
\bodytext{En la ejecución de prueba con 1193 alertas reales:}
\begin{itemize}
  \item CRÍTICA: 1179 (98.8%)
  \item INFORMATIVA: 13 (1.1%)
  \item IRRELEVANTE: 1 (0.1%)
\end{itemize}

% Evidencias
\sectiontitle{Evidencias}
\bodytext{Archivos generados y artefactos relevantes:}
\begin{itemize}
  \item `data/processed/alertas_clasificadas_demo_*.json` (salida del pipeline)
  \item `src/reasoning_agent_lite.py`, `src/reasoning_agent.py` (full/optional)
  \item `web/index.html` actualizado con panel de flujo y `web_server.py` para demo local
  \item Logs de ejecución: `sprint3_reasoning_demo.py` (registro de procesamiento)
\end{itemize}

% Incidencias y Riesgos
\sectiontitle{Incidencias y Riesgos}
\bodytext{Limitaciones detectadas durante el sprint:}
\begin{itemize}
  \item El clasificador inicial está entrenado con un set reducido (18 ejemplos); requiere re-entrenamiento con datos validados para mejorar balance y reducir bias.
  \item Dependencia futura en servicios de traducción y notificación (Sprint 4) que puede requerir claves/servicios externos.
  \item Cambios en estructuras HTML de fuentes externas pueden afectar scrapers (riesgo mitigado con tests y manejo de excepciones).
\end{itemize}

% Recomendaciones
\sectiontitle{Recomendaciones}
\bodytext{Para el Sprint 4 se recomienda:}
\begin{itemize}
  \item Implementar motor de traducción local/ligero o integrar solución cloud según disponibilidad.
  \item Diseñar e integrar un módulo de notificaciones (WhatsApp Business API + SMTP) con colas y reintentos.
  \item Construir dataset de entrenamiento con alertas validadas manualmente y reentrenar el clasificador.
  \item Conectar `web/index.html` con endpoint `/api/alerts/latest` y preparar pruebas de integración end-to-end.
\end{itemize}

\vspace{1cm}
\bodytext{Documento generado: 4 de Mayo de 2026.}

\end{document}
